\documentclass[11pt]{article}

\usepackage[margin=1in]{geometry}
\usepackage{amsmath}
\usepackage{amssymb}
\usepackage{booktabs}
\usepackage[hidelinks]{hyperref}

\title{Combat Relations and Working Observations}
\author{Computational Total War}
\date{\today}

\begin{document}

\maketitle

\section{Scope}

This file records observed combat relationships for \emph{Total War:
WARHAMMER III} and the abstractions we intend to use in computational
analysis.  Unless stated otherwise, these notes concern land combat under the
baseline described in \texttt{data/unit\_stats/README.md}: patch 8.1.1, ultra
unit scale, rank zero, and no campaign, terrain, fatigue, or temporary-effect
modifiers.

The equations below distinguish three things that are easy to conflate:
the probability that one attack lands, the number of entities able to attempt
attacks, and the damage produced by a successful attack.

\section{Melee hit probability}

For a valid melee attack, the working hit-probability relation is
\begin{equation}
  p_{\mathrm{hit}}
  = \operatorname{clamp}\!\left(
      0.08,
      0.90,
      0.35 + \frac{MA_{\mathrm{eff}}-MD_{\mathrm{eff}}}{100}
    \right),
  \label{eq:hit-probability}
\end{equation}
where $MA_{\mathrm{eff}}$ and $MD_{\mathrm{eff}}$ are measured in the
percentage-point units shown on the unit card.  Equivalently, the unclamped
percentage is $35 + MA_{\mathrm{eff}}-MD_{\mathrm{eff}}$.

Relevant attack bonuses are included before the roll.  In particular, bonus
versus infantry applies against small entities and bonus versus large applies
against large entities.  These bonuses increase both effective melee attack
and weapon damage.  Charge bonus also increases both quantities while it
decays.

Entity count does not appear in Equation~\ref{eq:hit-probability}.  Two units
with the same effective melee attack have the same probability per otherwise
identical attack, regardless of whether they contain 60 or 120 entities.

\section{Flanking}

Flanking has both a directional statistical effect and a geometric effect.
For the directional effect, the defending entity's effective melee defence is
\begin{equation}
  MD_{\mathrm{eff}} =
  \begin{cases}
    MD,     & \text{attack from the front},\\
    0.6 MD, & \text{attack from a side},\\
    0.3 MD, & \text{attack from the rear}.
  \end{cases}
  \label{eq:directional-md}
\end{equation}
Thus a side attack reduces melee defence by 40 percent and a rear attack by
70 percent.  These are multiplicative changes to melee defence, not direct
damage multipliers and not fixed percentage-point additions to hit chance.

The test is made per attacking and defending entity, using the defending
entity's current facing.  Surrounding a unit therefore does not impose one
permanent melee-defence penalty on every entity in that unit.  An entity that
turns to face a former rear attacker can defend against that attacker with its
front-facing value, although it may thereby expose itself to an attacker on
the opposite side.

As an example, an entity with 30 melee attack fighting an entity with 50 melee
defence has the following probabilities before other modifiers:
\begin{center}
  \begin{tabular}{lrr}
    \toprule
    Attack direction & Effective melee defence & Hit probability \\
    \midrule
    Front & 50 & 15\% \\
    Side  & 30 & 35\% \\
    Rear  & 15 & 50\% \\
    \bottomrule
  \end{tabular}
\end{center}

Flank and rear attacks also produce unit-level leadership penalties.  These
are distinct from the entity-level melee-defence adjustment and should be
modeled in a morale layer rather than as additional damage.

\section{Contact geometry and attack volume}

Only entities that obtain a valid engagement and execute attacks contribute
attack rolls.  A unit's total entity count is therefore an upper bound, not
the usual number of simultaneous attackers.  Rear ranks commonly contribute
no direct melee attacks while the front rank is engaged.

Let $E_d$ denote the number of actively attacking entities whose attacks fall
in direction class $d \in \{\mathrm{front},\mathrm{side},\mathrm{rear}\}$,
and let $T$ be the attack interval.  A first-order approximation to landed
hits per second is
\begin{equation}
  H \approx
  \sum_d \frac{E_d}{T} p_d,
  \label{eq:hit-throughput}
\end{equation}
subject to
\begin{equation}
  \sum_d E_d \leq N_{\mathrm{living}},
\end{equation}
where $p_d$ is obtained from Equations~\ref{eq:hit-probability}
and~\ref{eq:directional-md}.

Equation~\ref{eq:hit-throughput} is an approximation because attack interval
is not a complete animation and targeting simulation.  Pathing, displacement,
knockdown, weapon reach, collision boxes, and animation behavior can prevent
nominally available entities from attacking on schedule.

Flanking normally increases attack volume by opening another portion of the
defender's perimeter.  Its practical benefit can therefore arise twice:
\begin{enumerate}
  \item more attackers can make rolls because more contact surface is
        available; and
  \item side and rear rolls face reduced effective melee defence.
\end{enumerate}

\section{Entity count and entity size}

Entity count affects the ceiling on simultaneous attacks, the number of
potential targets, casualty progression, and overkill.  As entities die, a
multi-entity unit can lose attack throughput even when its remaining entities
retain unchanged card statistics.

Physical entity size affects contact geometry rather than directly changing
the base hit formula.  A large target generally presents enough perimeter for
more small attackers to contact it, while large attacking entities consume
more frontage and fewer may fit along a given line.  Actual results remain
dependent on collision radius, spacing, reach, animation, mass, and knockdown.

The small/large classification has a separate statistical effect: it selects
whether bonus versus infantry or bonus versus large applies.  Consequently,
size can change hit probability indirectly through $MA_{\mathrm{eff}}$, even
though physical radius is absent from Equation~\ref{eq:hit-probability}.

\section{Splash and collision attacks}

Some entities can select multiple targets with one splash attack.  Listed
weapon damage is divided among the selected splash targets; the maximum splash
target count is therefore not a direct multiplier on total weapon strength.
Target-size eligibility, splash zones, and animation behavior further govern
which entities can be selected.

Collision attacks are a separate attack mechanism used by some chariots and
large entities.  They can generate melee attack rolls when an entity collides
with enemies during a charge or attack animation and have their own cooldown
and target-list parameters.  They should not be inferred solely from ordinary
attack interval or maximum splash targets.

\section{Missile collision and shield blocking}

Missile attacks do not use Equation~\ref{eq:hit-probability}.  A projectile is
launched into the battle simulation and must physically intersect a target.
The probability of doing so is governed by the projectile and firing-unit
profiles---including calibration distance, calibration area, marksmanship,
spread, velocity, firing arc, and homing behavior---together with target size,
formation, motion, range, terrain, and intervening entities.  Melee attack and
melee defence do not determine whether a missile connects.

Let $p_{\mathrm{col}}$ denote the probability that a launched projectile
collides with an eligible target.  This probability is scenario-dependent and
should not be identified directly with the normalized \texttt{accuracy}
column.  Calibration and physical geometry can dominate the displayed
accuracy value.

After a collision, an eligible shield may block the projectile completely.
For an incoming direction in which the shield applies, let $s$ be the missile
block chance expressed as a fraction.  Then
\begin{equation}
  p_{\mathrm{unblocked}} = 1-s.
  \label{eq:missile-unblocked}
\end{equation}
Ordinary shields are directional, so firing into an unprotected flank or rear
can set $s$ to zero even though the target has a nonzero card value.  Projectile
types may also bypass ordinary shield blocking.  Special attributes, such as
Directional Shield and Ballistic Plating, alter the protected directions or
the classes of projectile that can be blocked.  Shield blocking is a separate
all-or-nothing test; it is not missile resistance and does not merely reduce
the damage of every projectile.

\section{Missile damage resolution}

For one unblocked direct hit, let $B$ and $P$ be the projectile's base and
armour-piercing damage after applicable attacker-side modifiers.  Armour
affects $B$ but not $P$.  For target armour $A$, draw
\begin{equation}
  R_A \sim \operatorname{Uniform}(0.5A,A)
\end{equation}
and cap the rolled reduction at 100 percent.  Damage before resistances is
therefore
\begin{equation}
  D_{\mathrm{pre-resist}}
  = B\left(1-\frac{\min(R_A,100)}{100}\right)+P.
  \label{eq:missile-armour}
\end{equation}
There is no minimum surviving base damage when the armour roll reaches 100.
For $A\leq100$, the expected value simplifies to
\begin{equation}
  \mathbb{E}[D_{\mathrm{pre-resist}}]
  = B(1-0.0075A)+P.
  \label{eq:missile-armour-expectation}
\end{equation}
The expectation for $A>100$ should be evaluated from
Equation~\ref{eq:missile-armour}, because part of the armour-roll interval is
capped.

Missile bonus versus infantry or large is selected using the target's
small/large classification when the projectile has the corresponding bonus.
Multi-projectile weapons create separate collision and damage opportunities
for each projectile.  Consequently,
\texttt{projectiles\_\allowbreak per\_\allowbreak shot},
\texttt{shots\_\allowbreak per\_\allowbreak volley}, and burst behavior must not be folded into the
damage of a single projectile.

After armour, applicable resistance percentages are summed into $Q$ and the
result is capped at 90 percent:
\begin{equation}
  D_{\mathrm{final}}
  = D_{\mathrm{pre-resist}}
    \left(1-\min(0.90,Q)\right).
  \label{eq:missile-resistance}
\end{equation}
For an ordinary physical missile, $Q$ can contain physical resistance,
missile resistance, ward save, and fire resistance or weakness when the
projectile is flaming.  Magical missiles ignore physical resistance.  Spell
missiles use spell resistance and missile resistance, ignore physical
resistance, and also use ward save and any applicable fire resistance or
weakness.  Fire weakness is negative and can offset positive resistances.
Flaming damage additionally applies the On Fire status, which halves incoming
healing for ten seconds.  Damage components are rounded when applied to the
target entity, and damage beyond the entity's remaining health is lost rather
than transferred to another entity.

As a simple example, the normalized High Elf Archer projectile has $B=16$ and
$P=3$.  Against $A=80$, Equation~\ref{eq:missile-armour-expectation} gives
\begin{equation}
  16(1-0.60)+3=9.4
\end{equation}
expected damage for an unblocked direct hit before resistances.  If an
applicable shield has $s=0.55$, the conditional contribution per colliding
projectile is $0.45\times9.4=4.23$ before resistances.

\section{Missile explosions and penetration}

An explosive projectile can carry both direct projectile damage and a
separate explosion profile with its own base damage, armour-piercing damage,
and radius.  Direct-hit and explosion damage must be resolved as distinct
damage events against each affected entity.  The number of entities reached
by an explosion is geometric and must not be represented by multiplying its
listed damage by the nominal entity count of the target unit.

Penetrating projectiles may collide with multiple entities along their path,
and articulated or multi-part targets can present more than one collision.
These cases require projectile-level simulation or an explicitly calibrated
expected-hit count.  They cannot be represented faithfully by the ordinary
single-collision model alone.

\section{Working melee damage abstraction}

For ordinary non-splash attacks, a useful top-level decomposition is
\begin{equation}
  \mathrm{DPS}
  \approx
  \underbrace{E}_{\text{contact geometry}}
  \times
  \underbrace{\frac{1}{T}}_{\text{attack rate}}
  \times
  \underbrace{p_{\mathrm{hit}}}_{\text{accuracy}}
  \times
  \underbrace{\mathbb{E}[D\mid\mathrm{hit}]}_{\text{damage resolution}}.
  \label{eq:dps-decomposition}
\end{equation}
This decomposition should be evaluated separately for front, side, and rear
attackers when modeling flanking.  Splash and collision profiles require
separate event-level treatment rather than multiplying
Equation~\ref{eq:dps-decomposition} by the number of potential targets.

\section{Working missile damage abstraction}

Let $N_f$ be the number of entities currently able to fire, $n_c$ the number
of projectiles launched by one such entity during a complete firing cycle,
and $T_c$ the duration of that cycle.  A first-order direct-fire approximation is
\begin{equation}
  \mathrm{DPS}_{\mathrm{missile}}
  \approx
  \underbrace{N_f\frac{n_c}{T_c}}_{\text{projectile launch rate}}
  \times
  \underbrace{p_{\mathrm{col}}}_{\text{physical collision}}
  \times
  \underbrace{p_{\mathrm{unblocked}}}_{\text{shield test}}
  \times
  \underbrace{\mathbb{E}[D_{\mathrm{final}}\mid
                         \mathrm{unblocked\ hit}]}_{\text{damage resolution}}.
  \label{eq:missile-dps}
\end{equation}
This is a scenario-level abstraction rather than a unit-card identity.
Formation and line of sight can make $N_f$ smaller than the living entity
count.  The complete cycle includes animation, burst timing, and effective
reload; the nominal reload field alone does not establish $T_c$.  The count
$n_c$ must use the same cycle boundary: projectile, shot, volley, and burst
counts must not be multiplied together without resolving their nesting.
The shield and damage factors must be conditioned on the same collision
population, including its incoming directions and eligible targets.
The probability $p_{\mathrm{col}}$ varies with the target and engagement.  The
unit-card Missile Strength value is a rate-style summary and should not be
treated as damage per individual arrow or bullet.  Explosion, penetration,
contact effects, and multiple collision opportunities must be added as
separate event-level terms.

\section{Provisional native reconstruction of missile dispersion}
\label{sec:native-dispersion}

\subsection{Evidence and scope}

A static inspection on 17--18 September 2026 recovered the following arithmetic
and several caller branches.  This is a reconstruction, not an official
engine specification or a runtime-validated hit model.  Instruction-level
operations and constants are directly observed; the mappings to calibration
area, calibration distance, and marksmanship are corroborated by a generated
row reader matching the version-53 projectile export and by runtime constructor
copies.  The reader's registration has not been linked completely to a named
table.  Effective shooter accuracy remains an accessor interpretation.
Some modifiers and branch predicates
remain unidentified.  The executable identity and function addresses are
recorded in Section~\ref{sec:native-provenance}.  Steam's installed-build
manifest agrees with the repository baseline; this is installed-build
evidence, not a depot-level integrity attestation for the executable.

The route has two distinct randomization stages: a target-local aim-point
offset and angular perturbations of the launch vector.  A preceding
target-prediction branch can lead a moving target without in-flight homing.
Other callers,
projectile categories, and homing branches can add different behavior.
The category lookup identifies the special category-8 scatter branch as
rockets; the inspected arrow and Horror-bolt profiles use arrow (1) and axe
(4), respectively.  Both use the \texttt{dual\_low\_fixed} trajectory setting,
whose selector resolves to fixed or low according to a runtime state field.
These shared exclusions and settings do not establish identical complete
unit firing paths.

\subsection{Prelaunch target prediction}

In the inspected branch with a null homing-parameter pointer, the code
solves an initial trajectory, obtains its flight-time estimate, and advances
the target position by that time multiplied by target motion components.
Height/terrain adjustment, final aim offset, and a final trajectory solve
follow.  This is prelaunch target leading, not steering toward the target
after launch.  It does not establish that every target class or projectile
uses the same prediction routine.  The local jitter below should therefore
be understood relative to the branch's chosen or predicted target point,
not necessarily the target's current position.

\subsection{Target-local aim-point offset}

For the inspected eligible-target branch, the recovered local offsets are
\begin{align}
  \delta_x &= (2U_1-1)\,mJR,\\
  \delta_y &= (U_2-\tfrac12)\,mH.
  \label{eq:target-aim-offset}
\end{align}
Here the extracted aiming-size modifier is $m=0.9$.  The target scalars
$R$, $H$, and $J$ are interpreted as radius, height, and an intersection-radius
ratio; their full constructor mappings remain unverified.  The helper applies
$\delta_x$ horizontally perpendicular to the source-to-target bearing and
$\delta_y$ vertically.  This is a target-facing rectangular offset, not a
uniform ground-plane disk.

The observed pseudorandom generator updates an unsigned 32-bit state $z$ by
\begin{equation}
  z'=(214013z+2531011)\bmod 2^{32},
  \qquad U=\frac{\lfloor z'/2^{16}\rfloor}{65535}.
\end{equation}
Continuous uniform, independent $U_i$ are an analytical approximation to
these successive discrete draws, not an additional engine finding.
The target branch requires an existing target with type value 4; its full
class mapping remains unresolved.  A firing-loop caller can reuse an aim
offset across a burst, so shot-to-shot independence must not be assumed.

\subsection{Precision factor and effective calibration scale}

Define $\operatorname{clamp}_{\ell}^{u}(x)=\min(u,\max(\ell,x))$.
For the inspected shooter-present branch and distance $r\geq0.01$, the
reconstructed precision factor is
\begin{align}
  G(\theta)&=1-\frac{\operatorname{clamp}_{-90}^{90}
                         (\theta_{\mathrm{deg}})}{180},\\
  F(r)&=\operatorname{clamp}_{0.01}^{1}\!\left[
    M G(\theta)\frac{\max(0,a+b)}{100}
    \left(\frac{dV}{r}\right)^2-\pi_{\mathrm{pen}}
  \right].
  \label{eq:missile-precision}
\end{align}
Here $r$ is the distance scalar consistent with horizontal target distance;
$d$ is interpreted as calibration distance; $a$ and $b$ are interpreted as
effective shooter accuracy and projectile marksmanship.  The angle $\theta$
is the stored launch elevation.  The symbol $M$ collects two multiplicative
accessor results, $V$ is another unit modifier entering the distance ratio,
and $\pi_{\mathrm{pen}}$ is a conditional penalty.  Their complete semantic
mapping and live values remain unresolved.  For $r<0.01$, the inspected
function returns $F=1$ before this calculation.  When the shooter pointer
is null, the caller does not invoke this function: it retains a separate
key-value default previously associated with the tower-accuracy parameter.
Equation~\ref{eq:missile-precision} must not be extended to that branch.
In particular, $F$ is
\emph{not} a hit probability despite its numerical bounds.

For the record scalar $C$ interpreted as calibration area, an unsigned
firing-history counter $n(t)$ produces
\begin{equation}
  C_{\mathrm{eff}}(t)=
  \begin{cases}
    C, & n(t)=0,\\
    \max\{C[1-0.1n(t)],\,0.5C\}, & n(t)>0.
  \end{cases}
\end{equation}
The counter is copied from a shooter field that is incremented in a
firing/reload completion branch and cleared by separate state-management
routines.  Its exact event granularity and reset conditions remain unresolved:
it is not established as unit rank, burst index, or one increment per unit
volley.  A zero-initialized launch input does not establish a permanently zero
counter.  The subsequent arithmetic is
\begin{equation}
  S(r,t)=\frac{C_{\mathrm{eff}}(t)}{F(r)},\qquad L(r,t)=\sqrt{S(r,t)}.
  \label{eq:missile-calibration-scale}
\end{equation}
The quantity $S$ is an internal area-like dispersion scale, not a measured
ground-impact area.

\subsection{Angular perturbation and the quadratic regime}

Let $k$ be the aspect scalar, $s_\alpha$ the caller-supplied scatter scale,
and $D$ the native denominator.  The recovered widths are
\begin{align}
  \alpha_\psi&=2\arctan\!\left(\frac{L}{2kD}\right),&
  \alpha_\theta&=2\arctan\!\left(\frac{kL}{2D}\right),\\
  \Delta\psi&=(U_3-\tfrac12)\alpha_\psi s_\alpha,&
  \Delta\theta&=(U_4-\tfrac12)\alpha_\theta s_\alpha.
  \label{eq:missile-angular-scatter}
\end{align}
The vector helper adds these perturbations to azimuth and elevation and
preserves the launch vector's magnitude.  The aspect scalar is stored as
$k=1$ in the inspected image, without a dynamic-write audit; the traced
batch caller supplies $s_\alpha=1$.  An entity firing caller instead reads
the scale from the shooter and can skip the immediate scatter helper.
Equation~\ref{eq:missile-angular-scatter} is conditional on that helper being
invoked.  The denominator is reconstructed as
$D=\max(|v_y|,r)$.  A follow-up trace of the ordinary ballistic solution
places vertical launch velocity at state offset \texttt{+4} and horizontal
distance at \texttt{+0xa0}; the traced firing route reaches scatter without
overwriting that velocity vector.  This strengthens the operand mapping,
but does not explain the apparently mixed-unit comparison.  We retain the
native expression rather than silently replacing it by $r$.
The inverse-tangent identification is supported by
the helper's range reduction and ballistic uses, not a separate runtime test.

The traced entity setup has one branch that writes $s_\alpha=1$ and another
that writes zero or a settings-derived value according to the action state.
The skip flag starts cleared and subsequently receives an action-descriptor
bit; it also affects packed launch flags.  Neither constructor defaults nor
projectile category alone establish which settings a particular unit receives.
For finite ordinary inputs, zero scale removes this angular contribution,
not target-local jitter or all other branches.  Skipping one helper call is
likewise not a guarantee of perfect accuracy.

Within the shooter-present branch, with $M=V=1$, $\pi_{\mathrm{pen}}=0$,
$\theta=0$, $n(t)=0$, and $a+b>0$,
between the clamps in Equation~\ref{eq:missile-precision},
\begin{equation}
  S(r,t)=\frac{100C}{a+b}\left(\frac{r}{d}\right)^2.
  \label{eq:missile-quadratic-regime}
\end{equation}
Thus quadratic scaling has a native-code basis within a specified regime;
it is not an unrestricted power law.  For $C\geq0$, the clamps give
$C\leq S\leq100C$ under these assumptions.  Zero elevation is a simplifying
reference, not the exact ballistic solution for firing across flat terrain.

Neither this quadratic regime nor the angular perturbations establish a
uniform circular cone, a uniform impact disk, or a ratio of hit probabilities
equal to the inverse ratio of calibration areas.  Aim-point jitter,
trajectory, target spacing, and collision geometry remain separate.

\subsection{From launch dispersion to collisions and attrition}

For an ordinary single-collision projectile, a principled definition is
\begin{equation}
  p_{\mathrm{col}}(\mathcal S)=
  \mathbb E_{\omega\mid\mathcal S}
  \left[\mathbf 1\{
    \Gamma(\omega;\mathcal S)
    \text{ intersects an eligible target}
  \}\right],
  \label{eq:missile-collision-expectation}
\end{equation}
where $\mathcal S$ specifies the engagement, $\omega$ includes targeting and
launch randomization, and $\Gamma$ is the resulting flight path through the
moving scene.  This is a definition of the desired probability, not a claim
that the full targeting, flight, and collision algorithms have been recovered.
A simulation using approximate bodies and formation spacing can estimate
it without live telemetry, but its geometry remains a modeling assumption.
The shortcut $p_{\mathrm{col}}\approx\min(1,K/S)$, for an assumed coverage
parameter $K$, is only a surrogate and is not an observed engine rule.

For casualty progression, retain $N_f(t)$ explicitly in
Equation~\ref{eq:missile-dps}.  It is bounded by living entities but is not
automatically proportional to remaining unit health: wounded survivors and
dead models affect firing capacity differently.  Barrier absorption,
mitigation order, spillover, and regeneration require separate treatment.
This inspection does not validate treating barrier as additional ordinary
health or justify precise duel-completion times.

\section{Current data limitations}

The normalized schema currently contains entity count, attack interval,
small/large classification, mass, anti-infantry and anti-large bonuses,
maximum splash targets, missile base and armour-piercing damage, projectile
and volley counts, reload time, projectile velocity, spread, marksmanship,
calibration fields, damage tags, ammunition, and range.  The companion
projectile and explosion lookups additionally preserve collision and timing
fields, penetration and homing references, contact-effect references, and
explosion profiles.  A reference key is not a complete implementation of the
referenced behavior.  Entity dimensions and formation spacing can be
investigated in source tables but are not a complete collision model in the
normalized unit row.  Weapon reach, animation-event timing, detailed splash
and collision-attack behavior, and complete shield eligibility still require
further reconstruction.

The present data can therefore support scenario-based analysis in which
$E_{\mathrm{front}}$, $E_{\mathrm{side}}$, and $E_{\mathrm{rear}}$ are supplied
as assumptions.  It cannot yet derive those engagement counts exactly from
unit statistics alone.  For missiles it can calculate conditional damage and
nominal launch rate, but it cannot derive $N_f$, $p_{\mathrm{col}}$, explosion
coverage, or penetrating-hit counts exactly from the normalized row alone.
The native reconstruction in Section~\ref{sec:native-dispersion} narrows the
uncertainty in launch dispersion; it does not remove these limitations.

\section{Sources}

\subsection{Native-inspection provenance}
\label{sec:native-provenance}

Static inspection of \texttt{Warhammer3.exe}, 17--18 September 2026;
the executable hash was rechecked on 18 September.
The image base is \texttt{0x140000000}; addresses below are preferred-image
virtual addresses, not live relocated addresses.  Executable SHA-256
(concatenate the two lines):
\begin{quote}\ttfamily\small
b7315fa718fd84e2e018e2c4df06600e9\\
df0076156b474f148d9d558c939aa55
\end{quote}
Steam's \texttt{appmanifest\_1142710.acf} reports
both \texttt{buildid} and \texttt{TargetBuildID} as 24237342, with
\texttt{StateFlags} equal to 4, matching the repository's Steam-build baseline.
This does not independently attest the executable against depot hashes.
The arithmetic was inspected with Capstone 5.0.9 and pefile 2024.8.26 without
executing the binary or observing a live battle.

\begin{center}
\begin{tabular}{ll}
\toprule
Function or anchor & Preferred-image address\\
\midrule
Probable projectile-record constructor & \texttt{0x1404d75cc}\\
Named aiming-size parameter registration & \texttt{0x1405348db}\\
Target-local random offsets & \texttt{0x142f698d0}\\
Target-facing offset transform & \texttt{0x142f657a4}\\
Prediction and final-solve route & \texttt{0x142f63674}\\
Target-motion prediction helper & \texttt{0x142e4b4bc}\\
Precision-factor arithmetic & \texttt{0x142d1df1c}\\
Launch-elevation modifier & \texttt{0x142f64350}\\
Calibration-to-angular-width route & \texttt{0x142f654dc}\\
Angular random draws & \texttt{0x142f65300}\\
Direction perturbation helper & \texttt{0x142f62d14}\\
Batch firing route & \texttt{0x142f7c754}\\
Launch-input default constructor & \texttt{0x142f5ce2c}\\
\bottomrule
\end{tabular}
\end{center}

In the ballistic solution, \texttt{0x142f63eae} stores launch elevation,
\texttt{0x142f63ef5} stores the velocity-vector portion containing the vertical
component, and \texttt{0x142f63efc} stores horizontal distance.
The firing route reaches calibration/scatter at \texttt{0x142f7c956}.
The predictor uses flight time at state \texttt{+0xa4} and target motion
components at \texttt{+0xe0}, \texttt{+0xe4}, and \texttt{+0xe8}.
The area-reduction counter at state \texttt{+0x74} traces through
\texttt{0x142f5c990} and \texttt{0x142f5cd14} to launch-input \texttt{+0x4c},
initialized to zero by an eight-byte store at \texttt{+0x48}.
The entity launch builder \texttt{0x142f19f10} subsequently copies shooter
\texttt{+0xd54} into input \texttt{+0x4c}.  Entity firing/reload route
\texttt{0x142f4d588} increments that shooter field at \texttt{0x142f4dc94};
\texttt{0x14312d2df} and \texttt{0x14312d43b} clear it.  The reset-state names
and increment granularity remain unresolved.

The entity route calls scatter at \texttt{0x142f4db39} with scale from
shooter \texttt{+0xc88}, unless byte \texttt{+0xd04} skips the call.
Setter \texttt{0x142f32b0c} writes scale 1; \texttt{0x142f2b930} writes an
argument supplied by action update \texttt{0x142f48a04}.  That argument is
zero unless action field \texttt{+0x78} equals 6 and flag bit 0 at
\texttt{+0xa0} is set, in which case it reads settings field \texttt{+0x218}.
The choice between the two setters additionally depends on action flags and
a global-value comparison; their game-mode semantics remain unresolved.
Action kind 6 is not projectile category 6.
Function \texttt{0x142f438d0} copies action flag bit 3 into the shooter skip
byte, and \texttt{0x142f19f10} also uses that byte to alter launch flags.

The constructor copies source offsets \texttt{+0xa0} and \texttt{+0xa4} to
record offsets \texttt{+0x88} and \texttt{+0x8c}, corroborated as calibration
distance and area.  Source \texttt{+0x80} maps to marksmanship at record
\texttt{+0x68}.  Generated reader \texttt{0x14063973c}, registered through
thunk \texttt{0x140639734} by \texttt{0x1406f3b28}, matches these fields'
version-53 serialized order.  Category converter \texttt{0x1404d919c} uses
the table at \texttt{0x143bbf750}.  Trajectory selector
\texttt{0x142f6f980} resolves \texttt{dual\_low\_fixed} (3) to fixed (2)
when state \texttt{+0x78}>0, otherwise low (0).
The target-offset routine reads \texttt{+0x9c}, \texttt{+0xa0}, and
\texttt{[target+0x2d8]+0x64}.  The named aiming parameter is
\texttt{projectile\_aiming\_target\_size\_modifier}.
These pointers and offsets preserve the distinction between observed
arithmetic and inferred field names.  Schema-decoded projectile, land-unit,
battle-entity, unit-spacing, and key-value-rule records supplied the
table-side comparison; the RPFM schema descriptions are corroborating
metadata, not native engine source.

\subsection{Published sources for the preceding combat relations}

Creative Assembly, ``Feature Focus \#2: Damage (Part 1),'' especially the
entity attack roll, hit-chance limits, target-class bonuses, and flank/rear
melee-defence multipliers:
\href{https://community.creative-assembly.com/total-war/total-war-warhammer/blogs/6-feature-focus-2-damage-part-1}{Creative Assembly source page}.

Creative Assembly, ``Feature Focus \#1: Elevation,'' for the distinction
between units and their individual entities and for entity-level combat
calculations:
\href{https://community.creative-assembly.com/total-war/total-war-warhammer/blogs/11-feature-focus-1-elevation}{Creative Assembly source page}.

Creative Assembly, ``Total War: WARHAMMER III --- Patch 5.1.0,'' for splash
damage division and collision attacks:
\href{https://community.creative-assembly.com/total-war/total-war-warhammer/blogs/23-total-war-warhammer-iii-patch-5-1-0}{Creative Assembly source page}.

Creative Assembly, ``Total War: WARHAMMER III --- Update 5.0.0,'' for
projectile, volley, explosion, calibration, Directional Shield, and Ballistic
Plating examples:
\href{https://community.creative-assembly.com/total-war/total-war-warhammer/blogs/17-total-war-warhammer-iii-update-5-0-0}{Creative Assembly source page}.

Creative Assembly, ``Total War: WARHAMMER III --- Siege Proving Grounds,''
for physical projectile simulation and directional shield protection:
\href{https://community.creative-assembly.com/total-war/total-war-warhammer/blogs/76}{Creative Assembly source page}.

\end{document}
